\documentclass[acmsmall,screen,review,anonymous]{acmart}

\begin{document}

\title{Do the Marks Survive the Pipeline? (working title)}

\author{Anonymous Authors}

\begin{abstract}
[Scratch. One paragraph: AI-code watermarking exists, generated code goes through ordinary
dev-lifecycle operations (rename, format, lint, merge, package) before it ships, we don't know
if watermarks survive that. We test N schemes under a mutation battery and a lifecycle chain.
Headline finding goes here once the full-scheme run is done. Fill in last.]
\end{abstract}

\maketitle


\section{Introduction}

[Scratch outline for this section:
\begin{itemize}
    \item Motivation: LLMs generate code, need provenance signal, watermarking is one candidate.
    \item Gap: watermarks are evaluated on model output as-is; real code gets renamed, formatted,
      linted, squash-merged, packaged before anyone reads it again. Unclear if the mark survives that.
    \item The question this paper asks: does it survive an ordinary (non-adversarial) dev pipeline,
      not an attack.
    \item What we did, one sentence: N schemes, individual-operation battery + lifecycle chain,
      detection checked at every step.
    \item Note-to-self: mention the reserialization-vs-edit-content distinction here once results
      are locked, it's probably the headline. Don't commit to numbers yet.
    \item Contributions bullet list --- draft after Results is finished, not before.
\end{itemize}
]


\section{Background and Motivation}

[Scratch:
\begin{itemize}
    \item How LLM watermarking works generally (green-list bias / distortion-free sampling, detector
      statistic, threshold).
    \item Why code is a harder substrate than prose (grammar constraints, and code specifically gets
      put through automated tooling that prose doesn't).
    \item Why we're testing an ordinary pipeline rather than an adversarial removal attack --- this is
      the "weaker, prior condition" framing, needs one clean paragraph.
\end{itemize}
]


\section{Related Work}

[Scratch, cite as found:
\begin{itemize}
    \item Text watermarking schemes (green-list / distortion-free family) --- list the ones we
      implement against, one line each.
    \item Code-specific watermarking (entropy-gated, syntax-aware) --- same.
    \item Multi-bit / message-encoding code watermarks --- different problem, mention why excluded.
    \item Existing robustness evaluations --- note they're mostly adversarial-attack framed
      (paraphrase, translation), and that's the gap we're filling.
\end{itemize}
]


\section{Research Questions}

[Scratch --- these are close to final but check wording once Results is done:
\begin{itemize}
    \item \textbf{RQ1:} Individual ordinary operations --- does detectability change, does it vary by scheme?
    \item \textbf{RQ2:} Cumulative effect through a realistic sequence of operations, not just one at a time.
    \item \textbf{RQ3:} How much of any observed loss is the edit itself vs. an incidental
      implementation detail of how the edit was applied?
\end{itemize}
]


\section{Approach}

\subsection{System Design}
[Scratch: describe the harness --- scheme layer (common generate/detect interface over the
vendored implementations), mutation layer (list of operations, where the content comes from),
lifecycle driver (the step sequence), analysis layer (outcome classes, unit of analysis). Diagram
would help here.]

\subsection{Implementation}
[Scratch: which schemes, which vendored codebases, why they can't share a process, generation
model, host repository the prompts/corpus are drawn from, detection threshold. Table of schemes x
design family would go here.]


\section{Evaluation}

\subsection{Experimental Setup}
[Scratch: the two evaluation designs (independent-operation battery vs. lifecycle chain), sample
sizes per scheme, what's still running vs. finished.]

\subsection{Evaluation Metrics}
[Scratch: retention rate, never-broke vs. end-retained, first-break operation, signal-kept ratio,
unit-of-analysis note (distinct baseline vs. run).]


\section{Results}
\label{sec:results}

[Scratch --- not written yet. Subsections to fill in once the full run is done:
\begin{itemize}
    \item Which schemes produce a detectable baseline at all.
    \item Individual-operation battery, all schemes.
    \item Lifecycle chain, cumulative survival.
    \item The reserialization control (edit content vs. how the edit is applied).
\end{itemize}
Numbers go here once locked --- do not draft prose around unconfirmed results.]


\section{Discussion}

[Scratch --- depends entirely on Results. Likely threads to develop:
\begin{itemize}
    \item What actually drives survival --- signal strength at generation time vs. specific operations.
    \item Practical guidance for tooling, if the reserialization thread holds up.
    \item Anything methodological worth flagging (sample-size pitfalls, etc.) if it recurs.
\end{itemize}
]


\section{Threats to Validity}

\subsection{Internal Validity}
[Scratch: chain step order, sample independence, anything about the mutation/randomness setup.]

\subsection{External Validity}
[Scratch: generation model, host repository/language, parameter settings vs. published defaults.]

\subsection{Construct Validity}
[Scratch: how "survival" is operationalized, any confounds in how mutations are implemented.]

\subsection{Conclusion Validity}
[Scratch: anything about sample sizes, missing comparisons, unfinished scheme coverage.]


\section{Conclusion}

[Scratch --- write last, after Results and Discussion are locked.]


\bibliographystyle{ACM-Reference-Format}
\bibliography{references}

\end{document}
